%%%%%%%%%%  scan idx 61  |  folio 48  %%%%%%%%%%
non
seulement un groupe abélien, mais un $\lambda$-anneau, il s'ensuit par transport de structure que $\underline{K}(X)$ est aussi un $\lambda$-anneau, donc est muni d'une structure multiplicative et d'applications $\lambda^i$. Il importe de définir ces structures directement, sans passer aux fibrés vectoriels, ce qui sera fait ci-dessous.

Si $E$ est un fibré vectoriel sur $X$, on désigne par $\gamma(E)$ sa classe dans $\underline{K}(X)$~; si $F$ est un faisceau algébrique cohérent sur $X$, $\gamma(\Gamma)$ désigne de même sa classe dans $\underline{K}(X)$. Si $Y$ est une sous-variété irréductible de $X$, on pose $\gamma(Y)=\gamma(\cO_Y)$, où $\cO_Y$ désigne le faisceau des anneaux locaux de $Y$, considéré comme faisceau algébrique cohérent sur $X$. On définit par linéarité le symbole $\gamma(Y)$ lorsque $Y$ est un cycle quelconque sur $X$. Quand il y aura des risques de confusion, on écrira $\gamma_X$ pour $\gamma$.

Soit $X$ un espace algébrique quelconque. On introduit sur $\underline{K}(X)$ une \emph{filtration}, en désignant par $\underline{K}(X)^i$ le sous-groupe de $\underline{K}(X)$ engendré par les $\gamma(F)$ avec dim.supp.\ $F\leq n-i$ (où $n$ est la dimension de $X$). Il résulte du «~lemme de dévissage~» de Serre, sous la forme donnée dans [3], le résultat suivant~:

\textbf{Proposition 2.6.} \emph{Le groupe $\underline{K}(X)^i$ est engendré par les $\gamma(Y)$, où $Y$ parcourt l'ensemble des sous-variétés de dimension $\leq n-i$ de $X$. Si $X$ est irréductible, on a $G^0(\underline{K}(X))=\Z$.}

Ce dernier isomorphisme est donné par l'homomorphisme de $\underline{K}(X)$ dans $\Z$ défini à l'aide de la notion de \emph{rang} d'un faisceau algébrique cohérent.

Supposons maintenant $X$ \emph{non singulière}. S'inspirant des développements 
